% ============================================================
% settings.tex
% THIẾT LẬP CHUNG CHO BÁO CÁO
% ============================================================

% ------------------------------------------------------------
% KHỔ GIẤY VÀ LỀ
% ------------------------------------------------------------
\usepackage[
    a4paper,
    top=2.5cm,
    bottom=2.5cm,
    left=3cm,
    right=2cm
]{geometry}


% ------------------------------------------------------------
% FONT + TIẾNG VIỆT
% ------------------------------------------------------------
\usepackage{fontspec}
\setmainfont{TeX Gyre Termes} % Tương đương Times New Roman

\usepackage{polyglossia}
\setdefaultlanguage{vietnamese}
\setotherlanguage{english}   % cho phần tài liệu tham khảo APA


% ------------------------------------------------------------
% GIÃN DÒNG
% ------------------------------------------------------------
\usepackage{setspace}
\onehalfspacing


% ------------------------------------------------------------
% CÁC PACKAGE CẦN THIẾT
% ------------------------------------------------------------
\usepackage{graphicx}
\usepackage{tikz}
\usepackage{eso-pic}

\usepackage{xcolor}

\usepackage{booktabs}
\usepackage{tabularx}
\usepackage{array}
\usepackage{longtable}

\usepackage{enumitem}
\usepackage{titlesec}
\usepackage{tocloft}

\usepackage{caption}
\usepackage{float}

\usepackage{csquotes}
% \usepackage{microtype}   % tắt cho nhanh, tránh timeout Overleaf free

\usepackage{hyperref}
\usepackage{bookmark}


% ------------------------------------------------------------
% TÀI LIỆU THAM KHẢO APA 7
% ------------------------------------------------------------
\usepackage[
    backend=biber,
    style=apa,
    language=english,        % APA không có bản tiếng Việt
    sorting=nyt,
    natbib=true
]{biblatex}

\addbibresource{tailieu.bib}


% ------------------------------------------------------------
% MÀU
% ------------------------------------------------------------
\definecolor{UELBlue}{HTML}{1F5A94}
\definecolor{UELDarkBlue}{HTML}{174C7E}
\definecolor{GuideGray}{HTML}{6B7280}


% ------------------------------------------------------------
% HYPERLINK
% ------------------------------------------------------------
\hypersetup{
    colorlinks=true,
    linkcolor=black,
    citecolor=UELDarkBlue,
    urlcolor=UELDarkBlue,
    pdftitle={Do an cuoi ky - Phat trien Web Kinh doanh}
}


% ============================================================
% ĐỊNH DẠNG ĐOẠN VĂN
% ============================================================

\usepackage{indentfirst}          % thụt đầu dòng cả đoạn đầu tiên sau tiêu đề
\setlength{\parindent}{1.27cm}
\setlength{\parskip}{0pt}

\setlist[itemize]{
    leftmargin=1.2cm,
    itemsep=2pt,
    topsep=3pt
}

\setlist[enumerate]{
    leftmargin=1.2cm,
    itemsep=2pt,
    topsep=3pt
}


% ============================================================
% ĐÁNH SỐ MỤC
%
% I.
% II.
%   II.1.
%   II.2.
% III.
% ============================================================

\renewcommand{\thesection}{\arabic{section}}

\renewcommand{\thesubsection}{\thesection.\arabic{subsection}}

\renewcommand{\thesubsubsection}{\thesubsection.\arabic{subsubsection}}

% Cho phép đánh số đến subsubsection
\setcounter{secnumdepth}{3}


% ============================================================
% ĐỊNH DẠNG HEADING TRONG NỘI DUNG
% ============================================================

% SECTION
\titleformat{\section}
    {\bfseries\fontsize{14}{18}\selectfont}
    {\thesection.}
    {0.6em}
    {}

\titlespacing{\section}
    {0pt}
    {18pt}
    {8pt}


% SUBSECTION
\titleformat{\subsection}
    {\bfseries\fontsize{13}{17}\selectfont}
    {\thesubsection.}
    {0.6em}
    {}

\titlespacing{\subsection}
    {0pt}
    {12pt}
    {5pt}


% SUBSUBSECTION
\titleformat{\subsubsection}
    {\bfseries\fontsize{12}{16}\selectfont}
    {\thesubsubsection.}
    {0.6em}
    {}

\titlespacing{\subsubsection}
    {0pt}
    {10pt}
    {4pt}


% ============================================================
% MỤC LỤC
% STYLE GẦN GIỐNG TEMPLATE MẪU
% ============================================================

% polyglossia tự đặt lại tên theo tiếng Việt, nên phải ghi đè trong \captionsvietnamese
\gappto\captionsvietnamese{\renewcommand{\contentsname}{MỤC LỤC}}

% Chỉ hiện Section và Subsection
% Không cho 1.1.1 chen vào mục lục
\setcounter{tocdepth}{2}


% ------------------------------------------------------------
% TIÊU ĐỀ "MỤC LỤC"
% ------------------------------------------------------------

\setlength{\cftbeforetoctitleskip}{0pt}

\renewcommand{\cfttoctitlefont}{%
    \hfill\bfseries
    \fontsize{14}{18}\selectfont
}

\renewcommand{\cftaftertoctitle}{\hfill}   % căn giữa chữ MỤC LỤC
\setlength{\cftaftertoctitleskip}{0.9cm}


% ------------------------------------------------------------
% SECTION TRONG MỤC LỤC
%
% I.     MỞ ĐẦU ............................... 1
% II.    ĐỊNH HƯỚNG NGHỀ NGHIỆP ............. 2
% ------------------------------------------------------------

\renewcommand{\cftsecfont}{
    \bfseries
    \fontsize{13}{17}\selectfont
}

\renewcommand{\cftsecpagefont}{
    \bfseries
    \fontsize{13}{17}\selectfont
}

% Khoảng cách giữa các section
\setlength{\cftbeforesecskip}{0.48cm}

% Không thụt section
\setlength{\cftsecindent}{0cm}

% Khoảng dành cho I., II., III.
\setlength{\cftsecnumwidth}{1.7cm}

% Thêm dấu chấm sau số La Mã
\renewcommand{\cftsecaftersnum}{.}

% Dấu chấm nối đến số trang
\renewcommand{\cftsecleader}{
    \cftdotfill{\cftdotsep}
}


% ------------------------------------------------------------
% SUBSECTION TRONG MỤC LỤC
%
%     II.1.   Khảo sát các hướng nghề nghiệp .... 2
% ------------------------------------------------------------

\renewcommand{\cftsubsecfont}{
    \normalfont
    \fontsize{12}{16}\selectfont
}

\renewcommand{\cftsubsecpagefont}{
    \normalfont
    \fontsize{12}{16}\selectfont
}

% Khoảng cách giữa subsection
\setlength{\cftbeforesubsecskip}{0.16cm}

% Thụt vào
\setlength{\cftsubsecindent}{0.9cm}

% Khoảng dành cho II.1.
\setlength{\cftsubsecnumwidth}{1.9cm}

% Dấu "." cuối II.1.
\renewcommand{\cftsubsecaftersnum}{.}

% Dấu chấm nối
\renewcommand{\cftsubsecleader}{
    \cftdotfill{\cftdotsep}
}


% Khoảng cách giữa các dấu chấm
\renewcommand{\cftdotsep}{1.5}


% ============================================================
% BẢNG
% ============================================================

\renewcommand{\arraystretch}{1.25}

\newcolumntype{Y}{
    >{\raggedright\arraybackslash}X
}

\renewcommand{\tablename}{Bảng}
\renewcommand{\figurename}{Hình}


% ============================================================
% CAPTION
% ============================================================

\captionsetup{
    font=small,
    labelfont=bf,
    justification=centering
}


% ============================================================
% PLACEHOLDER HƯỚNG DẪN
% ============================================================

\newcommand{\placeholder}[1]{%
    {\color{GuideGray}\itshape [#1]}%
}


% ============================================================
% FOOTER
% Trang 1, Trang 2, ...
% MÀU XANH UEL
% ============================================================

\usepackage{fancyhdr}

\pagestyle{fancy}

\fancyhf{}

\fancyfoot[C]{%
    \textbf{Trang \thepage}%
}

\renewcommand{\headrulewidth}{0pt}
\renewcommand{\footrulewidth}{0pt}


% Kiểu plain cũng dùng footer giống nhau
\fancypagestyle{plain}{
    \fancyhf{}

    \fancyfoot[C]{%
        \textbf{Trang \thepage}%
    }

    \renewcommand{\headrulewidth}{0pt}
    \renewcommand{\footrulewidth}{0pt}
}